\cleardoublepage
\chapternonum{摘要}
随着人工智能技术的大力发展，其逐渐融入了社会生活的方方面面。目前的深度学习模型早已不再是孤立地运行了，而是嵌入由软件应用、模型框架和加速硬件协同构成的复杂 AI 系统中。这种层级演进虽然使得模型运行效率大大增加，也降低了开发成本，然而各层级供应链的复杂性和层级之间的依赖性也引入了大量新的安全风险。针对 AI 系统这些供应链安全问题，现有学术界和工业界进行了大量的研究，然而他们大多集中于供应链软件的漏洞挖掘和生态分析方面，缺乏整体视角的研究框架。为此，本文不再局限于单一组件的漏洞或者安全性分析，而是从系统视角针对 AI 系统的软件应用层、模型框架层和硬件加速层供应链展开全面分析，针对各层级独特的难题设计了创新性的分析方法，系统揭示它们可能存在的新型攻击面并构建相应的检测或防御手段。

在软件应用层，针对 AI 系统软件供应链安全问题，本文突破了传统依赖包粒度的分析局限，首次系统性的研究了 Python 生态中模块粒度的模块冲突供应链威胁。为了解决软件供应链生态依赖图庞大且复杂和模块路径提取困难的难题，本文提出 ModuleGuard 分析框架，其包括了新颖的基于 AST 解析和动态安装模拟的模块提取技术（InstSimulator）和环境语义感知优化的依赖图解析技术（EnvResolver），并实现了针对整个 PyPI 生态 43 万软件包和 GitHub 3,711 个大型项目的大规模冲突检测，不仅揭示了模块冲突风险在生态中的演化特征，还识别出了 108 个高星项目存在模块冲突风险，更进一步揭示了由于包管理缺陷导致的软件供应链遭受的新型模块替换攻击，可以利用同名模块诱导用户下载恶意软件包并执行任意代码。

在模型框架层，针对 AI 系统模型供应链的安全问题，本文发现了一种基于框架合法 API 能力滥用的新型 TensorAbuse 攻击威胁。不同于已有研究依赖于不安全序列化格式或者显示恶意代码注入的方式，TensorAbuse 攻击具备更高的隐蔽性，无法被现有工具检测。为了系统性分析 TensorAbuse 攻击的影响能力和范围，本文提出了一套 TensorAbuse 攻击分析框架，其包含了持久化 API 提取技术（PersistExt）启发式地提取框架 API 源码，以及 LLM 赋能的隐藏能力分析技术（CapAnalysis），用以解决框架代码实现逻辑复杂和抽象度高的难题。该框架成功识别了 TensorFlow 中 31 个有助于构建恶意逻辑的 API，并构建出 5 类攻击原语，实现了 4 类隐蔽的恶意模型，成功绕过了主流检测工具和开源模型库的检测，获得 Huntr 1650 美元漏洞奖励。针对 TensorAbuse 攻击，本文还开发了恶意模型检测工具实现对模型供应链的防护。

在硬件加速层，针对 AI 系统底层驱动和固件供应链的安全问题，本文首次实现了针对 NVIDIA GPU ASLR 安全机制的逆向分析工作。为了攻克 NVIDIA GPU 的闭源环境和安全机制不透明的难题，本文提出了新颖的 GPU ASLR 逆向框架，其包括了基于启发式标志检索的 GPU 虚拟内存语义恢复逆向技术（FlagProbe）和基于锚点内存和指针追踪的随机化地址快速搜集技术（AnchorTrace）。该框架不仅首次逆向了完整的 CUDA 程序 GPU 虚拟地址空间布局和内存段语义，更是揭示了 NVIDIA GPU 在 ASLR 机制实现上存在低熵性和粗粒度随机化的缺陷。更进一步地，本文发现 GPU 侧地址空间布局与 CPU 侧地址强相关，并基于此缺陷实现了新型的利用 GPU 侧栈溢出漏洞实现的跨处理器攻击，能够泄漏 CPU 侧关键内存的随机化地址信息，从而削弱 CPU 侧的防护机制。

总的来说本文工作构建了一个覆盖 AI 系统从上至下全栈的供应链安全研究，不仅提出了多种创新技术和分析框架，开源了大量工具，为促进现实世界中 AI 系统的安全性提升提供了理论基础和解决方案。

\bigskip
\noindent \textbf{关键词}：人工智能供应链安全、软件供应链攻击、模型供应链安全、GPU 安全

\cleardoublepage
\chapternonum{Abstract}
With the rapid development of artificial intelligence (AI) technology, it has gradually integrated into all aspects of social life. Current deep learning models are no longer standalone artifacts. Instead, they are embedded within complex AI systems composed of the software applications, model frameworks, and acceleration hardware. While this layered architecture significantly improves computational efficiency and reduces development costs, it also introduces new security risks due to the complexity of individual supply chains and the strong interdependencies across layers. Existing research on AI system supply chain security predominantly focuses on software vulnerabilities or ecosystem-level dependency analysis, lacking a holistic, system-oriented perspective. To this end, this thsis moves beyond these aspects. Instead, it conducts a comprehensive analysis of the supply chains in the software application layer, model framework layer, and hardware acceleration layer of AI systems from a system perspective. It designs innovative analytical methods to address the unique challenges of each layer, systematically reveals potential new attack surfaces, and constructs corresponding detection or defense mechanisms.

At the software application layer, addressing the security issues of the software supply chain in AI systems, this thsis breaks through the limitations of traditional analysis at the dependency package granularity. It is the first to systematically study the module conflict threats in the Python ecosystem supply chain. To solve the problems of the vast and complex ecological dependency graph of the software supply chain and the difficulty of extracting module paths, this thsis proposes the ModuleGuard analysis framework. It includes a novel module extraction technique based on AST parsing and dynamic installation simulation (InstSimulator) and an environment-semantic-aware optimized dependency graph resolution technique (EnvResolver). It also implements large-scale conflict detection on over 430,000 software packages in the entire PyPI ecosystem and 3,711 large projects on GitHub, which not only reveals the evolutionary characteristics of module conflict risks in the ecosystem but also identifies 108 high-starred projects with such risks. Furthermore, it reveals a new type of module replacement attack on the software supply chain caused by package management defects, which can exploit homonymous modules to induce users to download malicious packages and execute arbitrary code.

At the model framework layer, addressing the security issues of the model supply chain in AI systems, this thsis discovers a new threat of the TensorAbuse attack, which is based on the abuse of legitimate API capabilities of the framework. Unlike existing research that relies on insecure serialization formats or explicit malicious code injection, the TensorAbuse attack is more stealthy and cannot be detected by current tools. To systematically analyze the impact and scope of the TensorAbuse attack, this thsis proposes a TensorAbuse attack analysis framework. It includes a persistent API extraction technique (PersistExt) to heuristically extract framework API source code that can be stored in models, and an LLM-powered hidden capability analysis technique (CapAnalysis) to address the challenges of complex implementation logic and high abstraction in the framework implementation. This framework successfully identified 31 APIs in the TensorFlow that could help construct malicious logic, built 5 types of attack primitives, implemented 4 types of stealthy malicious models, successfully bypassed mainstream detection tools and open-source model hubs, and earned a \$1,650 bug bounty from Huntr. For the TensorAbuse attack, this thsis also developed a malicious model detection tool to protect the model supply chain.

At the hardware acceleration layer, addressing the security issues of the underlying drivers and firmware supply chain in AI systems, this thsis is the first to conduct reverse analysis work on the NVIDIA GPU Address Space Layout Randomization security mechanism. To overcome the challenges of the closed-source environment and opaque security mechanisms of NVIDIA GPUs, this thsis proposes a novel GPU ASLR reverse engineering framework. It includes a reverse technique for GPU virtual memory semantic recovery based on heuristic flag construction and searching (FlagProbe) and a rapid collection technique for randomized addresses based on anchor memory and pointer tracing (AnchorTrace). This framework not only reverse-engineered the complete GPU virtual address space layout and memory segment semantics of a CUDA program for the first time but also revealed the defects of low entropy and coarse-grained randomization in the implementation of the ASLR mechanism on NVIDIA GPUs. Furthermore, this thsis discovered that the GPU-side address space layout is strongly correlated with the CPU-side address. Based on this flaw, a new type of cross-processor attack was implemented. By exploiting a GPU-side stack overflow vulnerability, this attack can leak the randomized address information of critical memory on the CPU side, thereby weakening the protection mechanisms on the CPU side.

In summary, this thsis establishes a full-stack supply chain security research covering AI systems from top to bottom. It not only proposes various innovative technologies and analysis frameworks, but also open-sources numerous tools. It provides a theoretical foundation and solutions for enhancing the security of real-world AI systems.

\bigskip
\noindent \textbf{Keywords}：AI Supply Chain Security, Software Supply Chain Attacks, Model Supply Chain Security, GPU Security